% ============================================================
%  Static Geometry and the SI Projection
%  Why natural units require no fractal corrections
%  Chapter content EN
%  Sourced: Dok. 011, 012, 013, 014, 015, 041, 044, 116, 133, 193, 258
% ============================================================

\title{Dok.~261 — Static Geometry and the SI Projection:\\
Why Natural Units Require No Fractal Corrections}
\author{Johann Pascher}
\date{May 2026}
\maketitle

\begin{center}
\textit{Context: Dok.~011 (fine structure), Dok.~013/014 (SI \& natural units),
Dok.~116 (Koide), Dok.~133 (fractal correction)} \\[2pt]
\textit{FFGFT archive: DOI 10.5281/zenodo.20117635}
\end{center}

\tableofcontents
\newpage

%=============================================================
\section{Setting}
%=============================================================

The 2019 SI reform fixed four constants ($c$, $\hbar$, $e$, $k_B$) to exact values
and thereby turned most dimensionful quantities into mere unit definitions. What
remains as \emph{genuine} empirical input are the dimensionless constants -- above
all the fine-structure constant $\alpha$, whose value is still measured after the
reform ($\mu_0$, $\varepsilon_0$ and $\alpha$ remained measured quantities). FFGFT
derives these dimensionless inputs as well, from the single parameter
$\xipar = \tfrac{4}{3}\times 10^{-4}$.

This document collects the conversion machinery compactly
(Tables~\ref{tab:si-fixiert-en}--\ref{tab:xi-inputs-en}) and then establishes one
central point: \emph{as long as one stays in natural units and pure ratios, the
description is static and relative and requires no fractal corrections.} The
fractal corrections appear only upon projection to absolute SI values.

%=============================================================
\section{The constants fixed by the SI reform}
%=============================================================

\begin{table}[htbp]
  \centering
  \caption{The constants fixed by the 2019 SI reform (the ``machinery'').
           They carry no physics, only scale, and become~$1$ in natural units.}
  \label{tab:si-fixiert-en}
  \begin{tabular}{llll}
    \toprule
    Constant & Value (exact since 2019) & defines & natural units \\
    \midrule
    $c$     & $299\,792\,458$ m/s                          & metre   & $=1$ \\
    $\hbar$ & $1{.}054\,571\,817\ldots\times10^{-34}$ J\,s & kg      & $=1$ \\
    $e$     & $1{.}602\,176\,634\times10^{-19}$ C           & ampere  & charge unit \\
    $k_B$   & $1{.}380\,649\times10^{-23}$ J/K              & kelvin  & $=1$ \\
    \bottomrule
  \end{tabular}
\end{table}

%=============================================================
\section{Quantity conversion in natural units}
%=============================================================

\begin{table}[htbp]
  \centering
  \caption{Quantity conversion in natural units ($\hbar=c=k_B=1$): all
           dimensionful quantities become powers of energy.}
  \label{tab:nat-umrechnung-en}
  {\small
  \begin{tabular}{lccl}
    \toprule
    Quantity & SI dimension & $[\hbar=c=k_B=1]$ & back to SI \\
    \midrule
    Energy        & J        & $E$            & base \\
    Mass          & kg       & $E$            & $\times\,S_{\mathrm{T0}}$ \ (via $c^{-2}$) \\
    Momentum      & kg\,m/s  & $E$            & $\times\,c^{-1}$ \\
    Length        & m        & $E^{-1}$       & $\ell[\mathrm{fm}]=197{.}327/E[\mathrm{MeV}]$ \\
    Time          & s        & $E^{-1}$       & $t[\mathrm{s}]=6{.}582\times10^{-22}/E[\mathrm{MeV}]$ \\
    Temperature   & K        & $E$            & $T[\mathrm{K}]=E[\mathrm{eV}]/(8{.}617\times10^{-5})$ \\
    Charge (Gauss)& C        & dimensionless  & $e=\sqrt{\alpha}\approx0{.}0854$ \\
    \bottomrule
  \end{tabular}}

  \medskip
  \footnotesize
  The conversion to SI runs through the \emph{single} mass scale
  $S_{\mathrm{T0}} = 1\,\mathrm{MeV}/c^2 = 1{.}782662\times10^{-30}$~kg (Dok.~013/014);
  the length factor is $\hbar c/S_{\mathrm{T0}} = 1{.}973\times10^{-13}$~m, i.e.\ the
  $197$\,MeV\,fm of the $\hbar c=1$ column expressed in T0 language.\par\smallskip
  In T0 units the charge unit is chosen ($e=\sqrt{4\pi\varepsilon_0\hbar c}$) so that
  the \emph{written} value is $\alpha=1$ (Dok.~044) -- a convention, not a physical
  intervention. The physical invariant $\alpha\approx1/137$ is reconstructed from
  $\xipar$ (Table~\ref{tab:xi-inputs-en}).
\end{table}

%=============================================================
\section{The dimensionless inputs and their derivation from $\xipar$}
%=============================================================

\begin{table}[htbp]
  \centering
  \caption{The surviving dimensionless inputs of physics and their derivation from
           the single parameter $\xipar=\tfrac{4}{3}\times10^{-4}$. All relations
           are sourced.}
  \label{tab:xi-inputs-en}
  {\footnotesize\setlength{\tabcolsep}{4pt}
  \begin{tabular}{lccl}
    \toprule
    Invariant & Standard physics & FFGFT relation (from $\xipar$) & Source \\
    \midrule
    $\xipar$           & ---            & $=\tfrac{4}{3}\times10^{-4}=4/30000$ & Dok.~013 \\
    $\alpha$ (SI)      & $1/137{.}036$  & $\xipar\,(\Ezero/1\,\mathrm{MeV})^2$\,$^{\dagger}$ & Dok.~011 \\
    Lepton masses      & measured       & $m=r\,\xipar^{\,p}\,v$ & Dok.~116, 006 \\
    \quad exponents    &                & $(p_e,p_\mu,p_\tau)=(\tfrac{3}{2},\,1,\,\tfrac{2}{3})$ & Dok.~116 \\
    \quad ratios       &                & $(r_e,r_\mu,r_\tau)=(\tfrac{4}{3},\,\tfrac{16}{5},\,\tfrac{25}{9})$ & Dok.~116 \\
    Koide $Q$          & $0{.}66666$    & $=\tfrac{2}{3}$ \ (consequence of the $\xipar$ exponents) & Dok.~116, 258 \\
    $\sin^2\theta_W$   & $0{.}2312$     & $\dfrac{1-\sqrt{1-4\alpha_W}}{4}$,\ \ $\alpha_W=\sqrt{\xipar}$ & Dok.~041 \\[0.6em]
    $\theta_{13}$      & $8{.}57^\circ$ & $\arcsin\!\big(\xipar^{1/3}\big)$ & Dok.~041 \\
    $G$                & measured       & $\dfrac{\xipar^2}{4m_e}\,\Kfrak$,\ \ $\Kfrak=1-100\xipar\approx0{.}9867$ & Dok.~012, 133 \\
    \bottomrule
  \end{tabular}}

  \medskip
  \footnotesize
  $^{\dagger}$\,$\alpha=\xipar\,(\Ezero/1\,\mathrm{MeV})^2$ is the \emph{leading order},
  with $\Ezero=\sqrt{m_e\,m_\mu}=7{.}398$\,MeV.
  CODATA precision is reached through the fractal correction product
  $\displaystyle\prod_{n=1}^{137}\!\big(1+\delta_n\,\xipar\,(4/3)^{\,n-1}\big)$.
  These corrections are \textbf{not free assumptions}: $\Kfrak=1-100\xipar$ is uniquely
  fixed by the scale flow from the Planck to the electroweak scale
  ($D_f^{\,\mathrm{space}}=3-\xipar$, Dok.~133, ``$D_f$ is uniquely determined --
  not freely choosable''), and the rational ladder ratios $r_i$ are
  $\alpha$-geometric invariants (Dok.~258), structurally inaccessible to any
  after-the-fact fitting (non-closure theorem, Dok.~193).
\end{table}

%=============================================================
\section{Static geometry: why natural units are correction-free}
%=============================================================

As long as physics is expressed in natural units and pure ratios, its description
is \emph{static} and \emph{relational}: only dimensionless ratios occur, no
absolute anchors and no scale dependence. At this level the FFGFT quantities are
exact geometric invariants of $\xipar$ -- and no fractal correction is required.

\begin{principle}
\textbf{Principle.} In natural units FFGFT is purely relational. The dimensionless
structural relations -- mass ratios, mixing angles, the Koide quantity -- are
static geometric invariants of $\xipar$. They contain no fractal corrections.
\end{principle}

The clearest example is the Koide quantity: $Q=\tfrac{2}{3}$ follows exactly from
the exponents $(p_e,p_\mu,p_\tau)=(\tfrac{3}{2},1,\tfrac{2}{3})$ and the ratios
$(r_e,r_\mu,r_\tau)=(\tfrac{4}{3},\tfrac{16}{5},\tfrac{25}{9})$ -- with no
additional parameter and no fractal correction (Dok.~116). Likewise the ladder
ratios $r_i$, the exponents $p_i$, the Weinberg structure and $\theta_{13}$ are
exact functions of $\xipar$ alone. None of these relations needs an absolute
scale, an $\Ezero$ in MeV, or a conversion factor.

Where, then, do the fractal corrections arise? Only at the \emph{SI projection} --
where an absolute scale is fixed. $\Kfrak=1-100\xipar$ is not a correction to the
geometry but the cumulative effect of the scale flow from the Planck to the
electroweak scale (Dok.~133): a property of mapping the geometry onto an absolute
energy ruler. Likewise the fractal correction product
$\prod_n(1+\delta_n\,\xipar\,(4/3)^{n-1})$ appears only when the leading relation
$\alpha=\xipar(\Ezero/1\,\mathrm{MeV})^2$ is brought to CODATA precision in absolute
SI values.

An important delimitation: these projection factors are nonetheless \emph{not}
free parameters. $\Kfrak$ is uniquely fixed by the scale flow (Dok.~133), the
rational ratios are $\alpha$-geometric invariants (Dok.~258), and the non-closure
theorem (Dok.~193) blocks any after-the-fact adjustment. They are derived
projection factors, not assumed fits.

\begin{result}
\textbf{Conclusion.} The fractal corrections are the price of the absolute
projection, not a defect of the geometry. As long as one stays in natural units
and ratios, the description is static, relative and correction-free; the
corrections appear exclusively at the transition to absolute SI values -- and even
there they are uniquely determined by $\xipar$ and the scale structure.
\end{result}
